\chapter{Juízos, Babilônia e a queda dos ídolos}
\noindent\textit{Vinheta.} Em uma sexta-feira, uma grande plataforma de \emph{e-commerce} cai após uma sequência de falhas. Em horas, pequenos lojistas ficam sem caixa; transportadoras reprogramam coletas; propagandas param. Nos noticiários, especialistas dizem: ``\emph{foi só uma instabilidade}''. No bairro, o efeito é real: \textbf{ninguém vende}, \textbf{ninguém compra}. Lembra Ap~18: mercadores lamentando a queda repentina da grande cidade-comércio. Não estamos em Apocalipse 18 literalmente, mas o \emph{padrão} nos instruirá hoje.

\section{Mapa bíblico (o que é não negociável)}
\begin{itemize}[leftmargin=*,noitemsep]
  \item \textbf{Deus julga com justiça} (\emph{selos, trombetas, taças} — Ap~6--16): juízos \emph{reais}, não caprichos; expõem idolatria e chamam ao arrependimento (2~Pe~3:9; Rm~2:4--5).
  \item \textbf{Babilônia cai} (Ap~17--18): sistema que mistura \emph{poder, luxo, violência e idolatria}; sedução econômica e \textbf{exploração} — inclusive de \emph{corpos e “almas humanas”} (Ap~18:13).
  \item \textbf{O Cordeiro vence} (Ap~19): a queda dos ídolos abre espaço para \textbf{aliança e festa} (bodas do Cordeiro).
  \item \textbf{Nova criação} (Ap~21--22): juízo não é fim frio; é \textbf{porta} para o mundo vindouro.
\end{itemize}

\section{Quatro propósitos dos juízos (leitura pastoral)}
\begin{checklist}
\begin{itemize}[leftmargin=*,noitemsep]
  \item \textbf{Expor} ídolos (o que adoramos quando achamos que é “só economia”).
  \item \textbf{Convidar} ao arrependimento (interrupções que acordam a cidade).
  \item \textbf{Proteger} o povo (limitar maldade; encorajar perseverança).
  \item \textbf{Remover} o mal (preparar o terreno para a Nova Jerusalém).
\end{itemize}
\end{checklist}

\section{Mini-exegese 1 — A lista de mercadorias (Ap 18:11--13)}
O catálogo de Ap~18 vai do \textbf{luxo} (ouro, púrpura, marfim) ao \textbf{cotidiano} (trigo, gado, azeite) e encerra com um soco: ``\emph{e corpos e almas de homens}''. Babilônia trata gente como \textbf{recurso}. Sempre que a economia \emph{exige} sacrificar pessoas na engrenagem, estamos “cheirando” Babilônia.

\section{Mini-exegese 2 — ``Sai dela, povo meu'' (Ap 18:4)}
O chamado não é isolamento geográfico, mas \textbf{separação de lealdade} e \textbf{práticas}. Permanecemos na cidade (Jr~29:7), porém \emph{não} participamos de seus pecados: \textbf{honestidade}, \textbf{justiça} e \textbf{generosidade} marcam nossa vida econômica.

\begin{nicbox}
\textbf{NIC aplicado}\\
\textbf{Narrativa}: Babilônia promete “segurança e luxo para sempre”. Deus narra: \emph{em uma hora} cai o ídolo.\\
\textbf{Identidade}: não somos mercadoria; somos \textbf{selados} para Deus (Ap~7).\\
\textbf{Controle}: portas de compra e venda se tornam instrumentos de coerção — respondemos com \textbf{fidelidade e partilha}.
\end{nicbox}

\section{O que evitar (atalhos ruins de interpretação)}
\begin{checklist}[title={Sem sensacionalismo}]
\begin{itemize}[leftmargin=*,noitemsep]
  \item \textbf{Datismo}: marcar datas para juízos específicos.
  \item \textbf{Nomear vilões imediatos}: “Babilônia é \emph{esta} empresa ou \emph{este} país” — texto não licencia reduções simplistas.
  \item \textbf{Tecnologismo}: reduzir símbolos espirituais a um único artefato (cap.~5 e 8 já mostraram o equilíbrio).
\end{itemize}
\end{checklist}

\section{Ferramenta — Inventário de idolatria econômica (30 minutos)}
\begin{checklist}
\begin{itemize}[leftmargin=*,noitemsep]
  \item \textbf{Coração}: o que me dá segurança (saldo, status, “ser visto”)?
  \item \textbf{Tempo}: onde meu calendário sacrifica família/igreja por ganho vazio?
  \item \textbf{Dinheiro}: %
  qual gasto evidencia adoração (ostentação) em vez de mordomia?
  \item \textbf{Confissão}: escreva 1 idolatria concreta; compartilhe com alguém de confiança; ore pedindo liberdade.
\end{itemize}
\end{checklist}

\section{Ferramenta — Liturgia de orçamento: simplicidade e generosidade}
\begin{checklist}[title={Passos práticos (família/solo)}]
\begin{itemize}[leftmargin=*,noitemsep]
  \item \textbf{Primeiro dez} (10\% do tempo/dinheiro): separe logo no início para o Reino e para os pobres.
  \item \textbf{Três caixas}: Essencial (moradia/alimentação/saúde), Missão (igreja/serviço), Futuro (reserva/estudos).
  \item \textbf{Corte um luxo} por trimestre; troque por hospitalidade/serviço.
  \item \textbf{Margem de misericórdia}: reserve 1--2\% para socorros inesperados (cap.~14 e 11).
\end{itemize}
\end{checklist}

\section{Ferramenta — Auditoria de justiça (igreja e negócios)}
\begin{checklist}[title={Perguntas que doem (e curam)}]
\begin{itemize}[leftmargin=*,noitemsep]
  \item Pagamos \textbf{no prazo} e com \textbf{clareza} a quem nos serve?
  \item Nossos preços/comunicação têm \textbf{transparência} (sem letrinhas para enganar)?
  \item Há espaço formal para \textbf{objeção de consciência} na empresa (cap.~15)?
  \item A igreja mantém \textbf{prestação de contas} pública das doações?
\end{itemize}
\end{checklist}

\section{Perguntas para grupos pequenos}
\begin{itemize}[leftmargin=*,noitemsep]
  \item Onde você sente a pressão de “ser Babilônia” na sua rotina de trabalho/consumo?
  \item Que \textbf{pequena renúncia} (dinheiro/tempo) poderia libertar seu coração esta semana?
  \item Como a sua igreja pode ser \textbf{boa notícia econômica} no seu bairro?
\end{itemize}

\bigskip
\noindent\textit{Próximo capítulo: Mapa das posições escatológicas — convergências, diferenças e por que isso importa menos do que você pensa.}
